\documentclass[12pt]{article}
\usepackage{amsmath, amssymb, amsthm}
\usepackage{geometry}
\usepackage{hyperref}
\usepackage{enumitem}
\usepackage{mathtools}

\geometry{margin=1.2in}

\theoremstyle{definition}
\newtheorem{definition}{Definition}
\newtheorem{theorem}{Theorem}
\newtheorem{proposition}{Proposition}
\newtheorem{corollary}{Corollary}
\newtheorem{remark}{Remark}
\newtheorem{observation}{Empirical Observation}
\newtheorem{conjecture}{Conjecture}

\title{\textbf{Evolutionary Receptor Topology Intelligence:\\
A Mathematical Formalization}\\
\large Receptor-Topological Dynamical Systems}
\author{ABI Research Program}
\date{}

\begin{document}

\maketitle

\begin{abstract}
We introduce the \textit{Receptor-Topological Dynamical System} (RTDS), a formal structure
for Evolutionary Receptor Topology Intelligence (ERTI). An RTDS is distinguished from
classical dynamical systems by three properties: (1) the state space dimensionality is not
fixed but changes under evolutionary selection pressure; (2) the evaluation function is
intrinsic to the state representation rather than externally specified; and (3) transitions
are learned from experience rather than designed. We formalize the receptor, the receptor
topology, the state space, the transition function, the intrinsic evaluation, the thinking
substrate, and the evolutionary operator. We state key results including the reachability
bound on topology growth and empirically observed dependency ordering of receptor
emergence. We distinguish theorems (proved from definitions) from empirical observations
(supported by experimental data, labeled as such) and conjectures (stated without proof).
\end{abstract}

\tableofcontents
\newpage

%% ============================================================
\section{Primitive: The Receptor}

Let $n$ denote the dimensionality of the raw observation vector from the environment.
$n$ is fixed within a generation (determined by the organism's body plan) but may change
across generations if the body plan evolves.

\begin{definition}[Receptor]
A \textbf{receptor} $r$ is a pair $r = (\phi_r, \alpha_r)$ where:
\begin{itemize}
    \item $\phi_r : \mathbb{R}^n \to [0,\infty)$ is the \textbf{activation function} mapping
    a raw observation $\mathbf{x} \in \mathbb{R}^n$ to a firing strength. The
    \textbf{detection condition} is the superlevel set
    $C_r = \{\mathbf{x} \in \mathbb{R}^n : \phi_r(\mathbf{x}) > \theta_r\}$ for
    threshold $\theta_r > 0$.
    \item $\alpha_r \in \mathbb{R}^+$ is the \textbf{metabolic cost} of maintaining $r$.
\end{itemize}
A receptor fires when $\phi_r(\mathbf{x}) > \theta_r$.
\end{definition}

\begin{remark}
$C_r$ is derived from $\phi_r$ and $\theta_r$, not an independent primitive. The codomain
of $\phi_r$ is $[0,\infty)$ rather than $[0,1]$ because some receptor activations are
unbounded (e.g., pain summed across limbs). The activation function includes any
normalization the receptor performs.
\end{remark}

\begin{remark}
The receptor fires \emph{when its condition is met}, not when it is lacking. This
distinguishes the receptor from a deficit signal. Detection precedes processing; the
thought, movement, or inner speech produced is downstream of the firing, not upstream.
\end{remark}

\begin{definition}[Receptor Fitness]
A receptor $r$ is \textbf{fitness-positive} in environment $\mathcal{W}$ if the expected
survival advantage of having $r$ exceeds its metabolic cost $\alpha_r$:
\[
\mathbb{E}_{\mathcal{W}}[\Delta \text{fitness} \mid r \text{ active}] > \alpha_r
\]
A receptor can only persist under selection if it is fitness-positive. This is the
load-bearing condition: the concept the receptor detects must be embedded in the
environment's causal structure with survival consequences attached.
\end{definition}

%% ============================================================
\section{Receptor Topology}

\begin{definition}[Receptor Topology]
A \textbf{receptor topology} at generation $t$ is a set
\[
\mathcal{R}_t = \{r_1, r_2, \ldots, r_{k_t}\}
\]
of fitness-positive receptors active in the organism at generation $t$, where $k_t = |\mathcal{R}_t|$
is the \textbf{topological dimension} at generation $t$.
\end{definition}

\begin{remark}
$k_t$ is not fixed. It increases when new receptors emerge and decreases when receptors
become fitness-negative. This is the reshaping phenomenon: complexity does not simply
expand the topology, it reorganizes it. Crucially, $k_t$ is \emph{non-monotone} ---
receptors can be gained and lost across generations. Empirically, across 60 generations:
$k_t$ ranged from 72 to 119, with receptors both gained (54) and lost (17) between
generation 0 and generation 59.
\end{remark}

\begin{definition}[Dependency Partial Order]
The receptor set $\mathcal{R}$ carries a \textbf{dependency partial order} $\preceq$ where
$r_i \preceq r_j$ means receptor $r_i$ must be active before $r_j$ can emerge. This
defines a directed acyclic graph $G_\mathcal{R} = (\mathcal{R}, \preceq)$ called the
\textbf{genome dependency graph}.
\end{definition}

\begin{proposition}[Layered Emergence]
Receptors emerge in layers consistent with $G_\mathcal{R}$. Define the layer of $r$ as:
\[
\ell(r) = \max_{r' \prec r} \ell(r') + 1, \quad \ell(r) = 0 \text{ if } r \text{ has no predecessors}
\]
Layer 0 receptors (trunk) emerge in any structured environment. Higher-layer receptors
require both the environmental structure and the presence of their prerequisite receptors.
\end{proposition}

\begin{observation}[Depth Prerequisite]
In a 60-generation ERTI run, depth\_reached ($\ell = 3$) activated at generation 29,
exactly one generation after metacognition and conflation ($\ell = 2$) became active at
generations 27--28. This is consistent with layered emergence but constitutes a single
empirical observation ($n=1$), not a proven theorem. Whether the ordering is a necessary
prerequisite or an artifact of the test design is an open question.
\end{observation}

%% ============================================================
\section{State Space}

\begin{definition}[Receptor State Space]
Given receptor topology $\mathcal{R}_t$, the \textbf{receptor state space} at generation
$t$ is:
\[
\mathcal{S}_t = \prod_{r \in \mathcal{R}_t} [0, \infty)
\]
a $k_t$-dimensional positive cone where each dimension corresponds to the activation
strength of one receptor.
\end{definition}

\begin{remark}
$\mathcal{S}_t$ is open-dimensional: its dimensionality $k_t$ changes across generations.
This is the defining structural property of an RTDS and distinguishes it from all fixed
state-space formalisms including finite state machines, Markov decision processes with
fixed state spaces, and standard dynamical systems. The codomain of $\phi_r$ and the
state space are consistent: both use $[0, \infty)$.
\end{remark}

\begin{definition}[State]
A \textbf{state} $s \in \mathcal{S}_t$ is a vector of receptor activation strengths:
\[
s = (\phi_{r_1}(\mathbf{x}), \phi_{r_2}(\mathbf{x}), \ldots, \phi_{r_{k_t}}(\mathbf{x}))
\]
where $\mathbf{x} \in \mathbb{R}^n$ is the raw observation vector from the environment.
We write $s_i = \phi_{r_i}(\mathbf{x})$ for the $i$-th component.
\end{definition}

%% ============================================================
\section{Transition Function}

\begin{definition}[True Transition Function]
The \textbf{true transition function} $T : \mathcal{S}_t \times \mathcal{A} \to \Delta(\mathcal{S}_t)$
maps a state-action pair to a distribution over next states, where $\mathcal{A}$ is the
action space. $T$ is unknown to the organism.
\end{definition}

\begin{definition}[Learned Transition Model]
The organism learns an approximation of $T$ via the \textbf{mental model}, a database of
cause-effect records:
\[
\mathcal{M} = \{(\mathbf{e}_i, \Delta s_i, \tau_i, c_i)\}_{i=1}^{N}
\]
where $\mathbf{e}_i$ is a context embedding, $\Delta s_i \in \mathbb{R}^{k_t}$ is the
observed state change (which may have negative components --- activation can decrease),
$\tau_i$ is the time delay, and $c_i \in [0,1]$ is the certainty.

The \textbf{readout function} $\hat{T} : \mathcal{S}_t \times \mathcal{A} \to \mathbb{R}^{k_t} \times [0,1]$
converts $\mathcal{M}$ into a point prediction:
\[
\hat{T}(s, a) = \left(\sum_{i \in \text{top-}k} w_i \cdot \Delta s_i, \quad
\frac{1}{k}\sum_{i \in \text{top-}k} c_i\right)
\]
where the top-$k$ entries are selected by cosine similarity between the query embedding
and stored embeddings $\mathbf{e}_i$, and $w_i$ are similarity-weighted coefficients
summing to 1. $\hat{T}$ returns a predicted state change in $\mathbb{R}^{k_t}$ and an
average certainty, not a distribution. The approximation $\hat{T} \approx \mathbb{E}[T]$
holds in the sense that the similarity-weighted average converges to the conditional
expectation as the store grows.
\end{definition}

\begin{remark}
$\hat{T}$ is context-conditioned: the same action $a$ from the same raw observation
can produce different $\Delta s$ depending on the current hidden state of the environment,
because the context embedding captures information beyond the action hash. The
contextual\_signal\_interpretation receptor (empirically discovered at generation 0)
is precisely the organism learning that $T$ is context-dependent --- that pain-in-context
is different from pain-as-absolute-signal.
\end{remark}

%% ============================================================
\section{Intrinsic Evaluation Function}

\begin{definition}[Intrinsic Evaluation]
The \textbf{intrinsic evaluation function} $V : \mathcal{S}_t \to \mathbb{R}$ maps any
state to a value determined by the receptor activation components of that state:
\[
V(s) = g(s_1, s_2, \ldots, s_{k_t})
\]
where $g : [0,\infty)^{k_t} \to \mathbb{R}$ is a function of the activation vector.
In the simplest case, $g$ is linear: $V(s) = \sum_{i} w_i \cdot s_i$ with valence
weights $w_i \in \mathbb{R}$. In general, $g$ may be nonlinear to capture receptor
interactions --- for example, simultaneous high pain and high endorphin (conflict)
produces a qualitatively different evaluation than either alone, which a linear $g$
cannot represent.
\end{definition}

\begin{remark}
$V$ is intrinsic in the strict sense: it is a property of the state representation itself,
not an external signal provided by a designer. The function $g$ (whether linear weights
or nonlinear interactions) is not specified --- it emerges from selection pressure.
In the ERTI implementation, the ArbitrationHead provides the nonlinear component,
modulating receptor group weights based on the current state.
\end{remark}

\begin{remark}[Evaluation Generalization]
$V$ produces a value for any state $s \in \mathcal{S}_t$, including states in environments
the organism has never encountered, because $V$ is defined over the activation vector
$s$ and receptors fire on conditions wherever those conditions are met. However, whether
that value is \emph{useful} --- whether acting on it leads to fitness --- depends on
whether the receptor topology $\mathcal{R}_t$ detects the structure that matters in the
new environment. Generalization of $V$ is necessary but not sufficient for
cross-environment transfer.
\end{remark}

%% ============================================================
\section{The Thinking Substrate}

\begin{definition}[Thinking Substrate]
The \textbf{thinking substrate} $M$ is a Monte Carlo Tree Search process over $(\mathcal{S}_t, \hat{T}, V)$:
\[
M(s) = \arg\max_{a \in \mathcal{A}} Q(s, a)
\]
where $Q(s,a)$ is estimated via UCB across simulated rollouts:
\[
Q(s, a) = \bar{V}(s, a) + c \sqrt{\frac{\ln N(s)}{N(s,a)}}
\]
with $\bar{V}(s,a)$ the mean value of rollouts from $(s,a)$, $N(s)$ the visit count of
state $s$, $N(s,a)$ the visit count of action $a$ from $s$, and $c$ an exploration
constant.
\end{definition}

\begin{definition}[Thinking Channels]
The MCTS tree generates \textbf{thinking channels} --- metadata about the search process
that become receptor inputs for the next cycle:
\begin{align*}
\chi_1 &= \max_{a} \bar{V}(s,a) && \text{(best\_value)} \\
\chi_2 &= H(\{N(s,a)/N(s)\}_a) / \ln|\mathcal{A}| && \text{(visit\_entropy, normalized to $[0,1]$)} \\
\chi_3 &= \max\!\left(0,\; 1 - \text{std}(\{\bar{V}(s,a)\}_a)\right) && \text{(value\_convergence, clipped to $[0,1]$)} \\
\chi_4 &= \max_a \bar{V}(s,a) - \min_a \bar{V}(s,a) && \text{(path\_divergence)} \\
\chi_5 &= |\{a : N(s,a) < 2\}| / |\mathcal{A}| && \text{(underexplored)} \\
\chi_6 &= d_{\max} / d_{\text{budget}} && \text{(depth\_reached)}
\end{align*}
All channels are clipped to $[-1, 1]$ before entering the observation vector.
These six channels are appended to $\mathbf{x}$, making the organism's thinking
a receptor input on the next cycle.
\end{definition}

\begin{remark}[Relationship to AlphaZero]
Standard MCTS (e.g., AlphaZero) requires an externally specified evaluation function
approximated by a value network trained on labeled outcomes. In ERTI, $V$ is intrinsic ---
no value network is needed because the receptor topology provides the evaluation.
Different receptor topologies evaluate identical states differently, and as the topology
develops across generations, the search improves because the organism gets better at
knowing what is worth searching for.
\end{remark}

%% ============================================================
\section{The Evolutionary Operator}

\begin{definition}[Evolutionary Operator]
The \textbf{evolutionary operator} $\mathcal{U}$ updates the receptor topology across
generations:
\[
\mathcal{U} : \mathcal{R}_t \times \mathcal{F}_t \to \mathcal{R}_{t+1}
\]
where $\mathcal{F}_t : 2^{\mathcal{R}^*} \to \mathbb{R}$ is the fitness function mapping
a receptor topology (a subset of the genome) to its observed fitness in environment
$\mathcal{W}_t$. The operator has three components:
\begin{enumerate}
    \item \textbf{Selection:} Fitness-proportional survival --- topologies with higher
    $\mathcal{F}_t$ are more likely to reproduce.
    \item \textbf{Topology bias inheritance:} Offspring inherit $\mathcal{R}_t$ as a
    prior, probe-gated: the inherited topology is validated against exploration data
    rather than accepted unconditionally.
    \item \textbf{Discovery:} Novel receptors emerge when the environment contains
    detectable structure not yet in $\mathcal{R}_t$ whose detection would be
    fitness-positive.
\end{enumerate}
\end{definition}

%% ============================================================
\section{The Full System: RTDS}

\begin{definition}[Receptor-Topological Dynamical System]
A \textbf{Receptor-Topological Dynamical System} (RTDS) is a tuple:
\[
\Sigma = (\mathcal{R}_0, \mathcal{G}, \hat{T}, V, M, \mathcal{U}, \mathcal{W}_t)
\]
where:
\begin{itemize}
    \item $\mathcal{R}_0$ is the initial receptor topology (trunk receptors)
    \item $\mathcal{G} = (\mathcal{R}^*, \preceq)$ is the genome --- the full dependency graph
    of all possible receptors ordered by their prerequisites
    \item $\hat{T}$ is the learned transition model (mental model with readout function)
    \item $V$ is the intrinsic evaluation function
    \item $M$ is the thinking substrate (MCTS over $\hat{T}$ and $V$)
    \item $\mathcal{U}$ is the evolutionary operator
    \item $\mathcal{W}_t$ is the environment at generation $t$
\end{itemize}
The system evolves across generations: $\mathcal{R}_{t+1} = \mathcal{U}(\mathcal{R}_t, \mathcal{F}_t)$,
and within each generation the organism operates as a dynamical system with state space
$\mathcal{S}_t$, transition model $\hat{T}$, evaluation $V$, and action selection $M$.
\end{definition}

%% ============================================================
\section{Key Results}

\begin{proposition}[Reachability]
Let $\Sigma$ be an RTDS. The set of receptors that can ever become active is bounded by
the genome:
\[
\forall t: \quad \mathcal{R}_t \subseteq \mathcal{R}^*
\]
For any receptor $r \in \mathcal{R}^*$, if there exists a generation $\tau$ such that
all prerequisites of $r$ are in $\mathcal{R}_\tau$ and $r$ is fitness-positive in
$\mathcal{W}_\tau$, then $r$ may enter $\mathcal{R}_{\tau+1}$.
\end{proposition}

\begin{remark}
Reachability is a direct consequence of how $\mathcal{U}$ is defined. It is stated
explicitly to clarify what the system can and cannot do: the genome $\mathcal{G}$
bounds the reachable topology. The ``ceiling'' of an RTDS is not
$|\mathcal{R}^*|$ (which is itself designer-specified and can grow --- e.g., the
contextual\_signal\_interpretation receptor was added to $\mathcal{R}^*$ after the
organism discovered it). The ceiling is bounded by the complexity of the environment
sequence $\{\mathcal{W}_t\}$. An RTDS with the environmental augmentation family
can increase its own environment's complexity, making the ceiling co-evolve with
the organism.
\end{remark}

\begin{remark}
Because $k_t$ is non-monotone (receptors can be lost when fitness-negative), the
sequence $\{k_t\}$ does not converge. The empirical record shows
$\limsup_{t \leq 60} k_t = 119$ and
$|\bigcup_{t \leq 60} \mathcal{R}_t| = 161$ out of $|\mathcal{R}^*| = 188$.
The cumulative union grows; the per-generation count oscillates.
\end{remark}

\begin{observation}[Topology Informativeness]
Different environment sequences produce different receptor topologies. Empirically,
social environments produce strong social receptor activation while tool-use
environments produce interaction-family activation. The topology is readable as a
partial record of what was worth detecting, though it is not a sufficient statistic
--- it does not capture which receptors were tried and lost, or the order of emergence.
\end{observation}

\begin{observation}[Metacognition Prerequisite for Depth]
In a 60-generation ERTI run, the depth\_reached thinking channel $\chi_6$ became
behaviorally influential at generation 29, one generation after metacognition and
conflation both activated (generations 27--28). This is consistent with the dependency
ordering in $G_\mathcal{R}$. This is a single empirical observation ($n=1$).
Whether the ordering is a necessary prerequisite or an artifact of the test design
is an open question.
\end{observation}

\begin{observation}[Reshaping Under Deep Time]
An RTDS operating under deep time learning produces qualitatively different cognitive
structures than single-run gradient-based training. Specifically: (a) the full
epistemic family (belief, doubt, conflation, epistemic strategy) activated under
evolutionary selection but not in single-run training; (b) 18 trunk receptors were
invariant across all 8 environment tiers; (c) the total number of simultaneously
active receptors per generation stabilized around 95--119 while the cumulative set
grew to 161. Complexity reshapes the topology rather than expanding it monotonically.
\end{observation}

\begin{observation}[Path-Dependent Topology]
The receptor topology is not uniquely determined by the environment. In a controlled
cross-environment transfer experiment, organisms evolved for 29 generations in
environment $\mathcal{W}_A$ were dropped into environment $\mathcal{W}_B$ alongside
naive organisms with no prior evolution. Both populations evolved for 10 generations
in the same $\mathcal{W}_B$. Result: 71 receptors discovered by both, 30 discovered
only by transfer organisms, 27 discovered only by naive organisms.

The trunk converged (71 shared). The canopy diverged based on evolutionary history.
Transfer organisms used their prior scaffolding to access deep canopy receptors
(composite affordance, hierarchical abstraction, niche construction) that naive
organisms couldn't reach. Naive organisms discovered receptors (metacognition,
theory of mind, epistemic strategy) that the transfer organisms' prior topology
biased them away from.

The topology is jointly determined by the environment \emph{and} the evolutionary
path. Different histories in the same environment produce genuinely different
but equally adequate intelligences. The trunk is universal. The canopy is biography.
\end{observation}

%% ============================================================
\section{The Three Relationships to the World}

The RTDS formalizes three orthogonal relationships an organism can have with its environment:

\begin{definition}[Exploration]
The organism moves through $\mathcal{W}_t$. The world does not change; $\hat{T}$ is
updated. State: $s \in \mathcal{S}_t$ changes; $\mathcal{W}_t$ does not.
\end{definition}

\begin{definition}[Interaction]
The organism acts on $\mathcal{W}_t$ and $\mathcal{W}_t$ responds. Both $s$ and
$\mathcal{W}_t$ change state.
\end{definition}

\begin{definition}[Environmental Augmentation]
The organism modifies $\mathcal{W}_t$ to increase its complexity, creating selection
pressure for receptors the organism does not yet have:
\[
\mathcal{W}_{t+1} = f(\mathcal{W}_t, \mathcal{R}_t)
\]
where $f$ increases the complexity of $\mathcal{W}$ in directions that will make new
receptors fitness-positive. This is the mechanism by which the organism builds the world
that builds the organism that builds the next world.
\end{definition}

\begin{remark}[Self-Augmentation]
A fourth relationship --- self-augmentation --- involves the organism changing its own
receptor topology or physical structure independently of environmental interaction.
In an RTDS, this corresponds to $\mathcal{R}_t$ changing within a generation rather than
only across generations via $\mathcal{U}$.
\end{remark}

%% ============================================================
\section{Cross-Domain Relevance Compression and Mathematics}

\begin{definition}[Cross-Domain Relevance Compression]
A receptor $r$ performs \textbf{cross-domain relevance compression} if its activation
function $\phi_r$ fires above threshold on states from multiple domains
$D_1, D_2, \ldots, D_m$ that share no surface features:
\[
\phi_r(\mathbf{x}_{D_1}) > \theta_r \wedge \phi_r(\mathbf{x}_{D_2}) > \theta_r \wedge \ldots
\]
where $\mathbf{x}_{D_j}$ is an observation in domain $D_j$. What makes the domains
``structurally equivalent with respect to $r$'' is defined by $\phi_r$ itself: two
domains are equivalent under $r$ if and only if $r$ fires in both. The receptor
compresses whatever invariant it detects across domains into a single activation.
\end{definition}

\begin{remark}
The mathematics receptor family implements cross-domain relevance compression at
increasing levels of abstraction. The quantity receptor fires on numerosity regardless
of what is being counted. The structural invariance receptor fires on relational patterns
regardless of the domain. The necessity receptor fires on the detection that a pattern
could not be otherwise --- proof structure is cross-domain relevance compression at its
deepest level. Mathematical ability requires cross-domain experience: the invariants
are only visible against sufficient domain variation.
\end{remark}

%% ============================================================
\section{The MCTS Bootstrapping Argument}

\begin{conjecture}[Self-Improving Loop]
In an RTDS with thinking substrate $M$, the following loop operates simultaneously
at multiple timescales:
\[
\mathcal{R}_t \to V_t \to M_t \to \text{richer experience} \to \hat{T}_{t+1}
\to \text{deeper thinking} \to \mathcal{R}_{t+1}
\]
Better receptor topology produces better intrinsic evaluation, which produces better
MCTS search, which produces richer experience and richer mental model, which enables
deeper thinking, which produces selection pressure for better receptor topology.
Whether this loop converges, diverges, or oscillates is an open question. The empirical
evidence from 60 generations shows oscillation in both $k_t$ and thinking influence
rather than monotonic convergence.
\end{conjecture}

\begin{remark}[What MCTS Contributes]
The cognitive prerequisites for the self-improving loop (metacognition, conflation,
fundamental distinction) emerge from environmental complexity \emph{without} MCTS ---
they appeared in deep time runs before the thinking substrate was added. What MCTS
contributes is making thinking quality \emph{visible to selection}: the six thinking
channels externalize the thinking process into the observation vector, allowing
organisms that think better to be selected over organisms that don't. MCTS does not
make the organism smarter. It makes the organism's smartness selectable.
\end{remark}

\begin{conjecture}[Receptor Topology as Sufficient Framing]
The self-improving loop requires an evaluation function that generalizes across every
level of the loop. The receptor topology provides such a function. Whether it is the
\emph{only} framing that can do so is an open question. A fixed reward function applied
uniformly at every level is a potential counterexample; whether such a function can
support the same developmental trajectory is untested. The claim is that the receptor
framing is \emph{sufficient}, not that it is uniquely necessary.
\end{conjecture}

%% ============================================================
\section{Relationship to Prior Work}

An RTDS shares components with several existing formalisms but is not reducible to any:

\begin{itemize}
    \item \textbf{Active inference} (Friston): shares intrinsic evaluation (free energy
    as self-generated objective) and learned world models. Differs in that active inference
    operates on a fixed state space; RTDS has open-dimensional $\mathcal{S}_t$. The
    receptor topology provides a finer-grained decomposition of the evaluation function
    than the single free energy functional.

    \item \textbf{NEAT and open-ended evolution} (Stanley \& Miikkulainen): shares
    growing network topology under evolutionary pressure. Differs in that NEAT grows
    the \emph{processing architecture} (hidden nodes and connections); RTDS grows the
    \emph{state representation} (what the organism can detect). The evaluation function
    in NEAT is externally specified; in RTDS it is intrinsic.

    \item \textbf{Nonparametric Bayesian models} (Dirichlet processes, Indian buffet
    processes): share data-driven dimensionality --- the model complexity grows with
    data. Differs in that nonparametric Bayes grows latent dimensions for statistical
    fit; RTDS grows receptor dimensions under \emph{fitness} pressure, which is a
    causal selection criterion rather than a statistical one.

    \item \textbf{MDPs and POMDPs}: share state-action-transition structure. Differ in
    that the state space, reward function, and transition function are all fixed and
    externally specified; in RTDS all three evolve.
\end{itemize}

The RTDS is distinguished by the \emph{conjunction} of open-dimensional state space,
intrinsic evaluation, and learned transitions. Each component exists in prior work;
the combination does not.

%% ============================================================
\section{Summary: What Makes RTDS Formally Distinct}

An RTDS is distinguished from prior formalisms by the conjunction of three properties:

\begin{enumerate}
    \item \textbf{Open-dimensional state space}: $k_t$ is not fixed. The dimensionality
    of the state space changes under evolutionary selection pressure according to the
    genome dependency graph $\mathcal{G}$. It can increase (new receptors) or decrease
    (receptors becoming fitness-negative).

    \item \textbf{Intrinsic evaluation}: $V$ is not externally specified. It emerges
    from the receptor topology and applies to any state in $\mathcal{S}_t$, though its
    usefulness in novel environments depends on whether $\mathcal{R}_t$ detects
    relevant structure.

    \item \textbf{Learned transition model}: $\hat{T}$ is not designed. It is built
    from the organism's experience via a readout function over stored cause-effect
    records, and updated continuously.
\end{enumerate}

\bigskip
\noindent\textit{Whoever has receptors, let them recept.}

\end{document}
